\newcommand{\mxA}{\mx{A}} % matrix A
\newcommand{\mxB}{\mx{B}} % matrix B
\newcommand{\mxC}{\mx{C}} % matrix C
\newcommand{\mxD}{\mx{D}} % matrix D
\newcommand{\mxE}{\mx{E}} % matrix E
\newcommand{\mxK}{\mx{K}} % matrix K
\newcommand{\mxL}{\mx{L}} % matrix L
\newcommand{\mxM}{\mx{M}} % matrix M
\newcommand{\mxN}{\mx{N}} % matrix N
\newcommand{\mxO}{\mx{O}} % matrix O
\newcommand{\mxQ}{\mx{Q}} % matrix Q
\newcommand{\mxP}{\mx{P}} % matrix P
\newcommand{\mxR}{\mx{R}} % matrix R
\newcommand{\mxS}{\mx{S}} % matrix S
\newcommand{\mxT}{\mx{T}} % matrix T
\newcommand{\mxW}{\mx{W}} % matrix W
\newcommand{\mxX}{\mx{X}} % matrix X
\newcommand{\mxY}{\mx{Y}} % matrix Y
\newcommand{\mxZ}{\mx{Z}} % matrix Z
\newcommand{\mxG}{\mx{G}} % matrix G
\newcommand{\mxH}{\mx{H}} % matrix H
\newcommand{\mxU}{\mx{U}} % matrix U
\newcommand{\mxV}{\mx{V}} % matrix V
\newcommand{\mxF}{\mx{F}} % matrix F
\newcommand{\mxSigma}{\mx{\Sigma}} % matrix Sigma
\newcommand{\mxOmega}{\mx{\Omega}} % matrix Omega

\newcommand{\vca}{\vc{a}} % vector a
\newcommand{\vcb}{\vc{b}} % vector b
\newcommand{\vce}{\vc{e}} % vector e
\newcommand{\vcs}{\vc{s}} % vector s
\newcommand{\vcu}{\vc{u}} % vector u
\newcommand{\vcv}{\vc{v}} % vector v
\newcommand{\vcw}{\vc{w}} % vector w
\newcommand{\vcx}{\vc{x}} % vector x
\newcommand{\vcy}{\vc{y}} % vector y
\newcommand{\vcz}{\vc{z}} % vector z
\newcommand{\vctheta}{\vc{\theta}} % vector \theta
\newcommand{\vcmu}{\vc{\mu}} % vector \mu
\newcommand{\vcp}{\vc{p}} % vector p
\newcommand{\vcq}{\vc{q}} % vector q

\newcommand{\Nu}{n_\vc{u}} % size of vector u
\newcommand{\Nv}{n_\vc{u}} % size of vector v
\newcommand{\Nw}{n_\vc{w}} % size of vector w
\newcommand{\Nx}{n_\vc{x}} % size of vector x
\newcommand{\Ny}{n_\vc{y}} % size of vector y


\newcommand{\hatA}{\hat{\mx{A}}} % hat of (matrix) A
\newcommand{\hatP}{\hat{\mx{P}}} % hat of (matrix) P
\newcommand{\hatQ}{\hat{\mx{Q}}} % hat of (matrix) Q
\newcommand{\hatX}{\hat{\mx{X}}} % hat of (matrix) X

\newcommand{\uGamma}{\underline{\Gamma}} % underline (matrix) Gamma
\newcommand{\bGamma}{\overline{\Gamma}} % overline (matrix) Gamma

%\newcommand{\nD}{{n_{\Dcal}}} % size of data set
\newcommand{\nD}{{n_{\smtiny{$\Dcal$}}}} % size of data set
\newcommand{\nsD}{{n_{\stiny{$\Dcal$}}}} % size of data set
\newcommand{\nmD}{{n_{\mtiny{$\Dcal$}}}} % size of data set
\newcommand{\nsmD}{{n_{\smtiny{$\Dcal$}}}} % size of data set
% \Dcal --> data set

\newcommand{\nbatch}{n_{\text{batch}}}

\newcommand{\mux}{\mu_{\vcx_0}}     % vector mu_x0
\newcommand{\mxSx}{\mx{S}_{\vcx_0}}   % vector S_x0
\newcommand{\mxSw}{\mx{S}_{\vcw}}     % vector S_w
\newcommand{\mxSv}{\mx{S}_{\vcv}}     % vector S_v


\newcommand{\thetaML}{\theta^{\text{\mtiny{$\mathrm{(ML)}$}}}}     % thetaML
\newcommand{\vcthetaML}{\theta^{\text{\mtiny{$\mathrm{(ML)}$}}}}     % thetaML

\newcommand{\Sbfxtheta}{\mxS^{{(\mathbf{x})}}\!(\theta)}
\newcommand{\Sbfytheta}{\mxS^{{(\mathbf{y})}}\!(\theta)}
\newcommand{\mubfxtheta}{\mu^{{(\mathbf{x})}}\!(\theta)}
\newcommand{\mubfythetax}{\mu^{{(\mathbf{y})}}\!(\theta, \mathbf{x})}
\newcommand{\Xitheta}{\Xi(\theta)}
\newcommand{\Dtheta}{\mathbf{D}(\theta)}
\newcommand{\Sqtheta}{\mxS_\vcq(\theta)}
\newcommand{\qtheta}{\vcq(\theta)}


\newcommand{\extendedR}{\overline{\Rbb}}

\newcommand\indep{\protect\mathpalette{\protect\indeP}{\perp}}
\def\indeP#1#2{\mathrel{\rlap{$#1#2$}\mkern2mu{#1#2}}}

\newcommand{\define}{\coloneqq}      % Definition: colon equality 
% \newcommand{\define}{\triangleq}      % Definition: triangle equality
\newcommand{\trans}[1]{{#1}^{\tr}}   % Transpose



% Style definitions for the different amsthm environments
\crefformat{equation}{#2(#1)#3}
\Crefformat{equation}{#2(#1)#3}
\crefrangeformat{equation}{#3(#1)#4-#5(#2)#6}
\Crefrangeformat{equation}{#3(#1)#4-#5(#2)#6}

% \theoremstyle{definition}
% \newtheorem{definition}{Definition}[section]
\crefformat{definition}{#2definition~#1#3}
\Crefformat{definition}{#2Definition~#1#3}
\crefrangeformat{definition}{#3definition~#1#4-#5#2#6}
\Crefrangeformat{definition}{#3Definition~#1#4-#5#2#6}

% \newtheorem{proposition}{Proposition}
\crefformat{proposition}{#2Proposition~#1#3}
\Crefformat{proposition}{#2Proposition~#1#3}
\crefrangeformat{proposition}{#3Proposition~#1#4-#5#2#6}
\Crefrangeformat{proposition}{#3Proposition~#1#4-#5#2#6}

% \newtheorem{theorem}{Theorem}[section]
\crefformat{theorem}{#2Theorem~#1#3}
\Crefformat{theorem}{#2Theorem~#1#3}
\crefrangeformat{theorem}{#3Theorem~#1#4-#5#2#6}
\Crefrangeformat{theorem}{#3Theorem~#1#4-#5#2#6}

% \newtheorem{lemma}{Lemma}
\crefformat{lemma}{#2Lemma~#1#3}
\Crefformat{lemma}{#2Lemma~#1#3}
\crefrangeformat{lemma}{#3Lemma~#1#4-#5#2#6}
\Crefrangeformat{lemma}{#3Lemma~#1#4-#5#2#6}

% \theoremstyle{remark}
% \newtheorem*{remark}{Remark}
\crefformat{remark}{#2Remark~#1#3}
\Crefformat{remark}{#2Remark~#1#3}
\crefrangeformat{remark}{#3Remark~#1#4-#5#2#6}
\Crefrangeformat{remark}{#3Remark~#1#4-#5#2#6}

% \newtheorem{assumption}{Assumption}
\crefformat{assumption}{#2Assumption~#1#3}
\Crefformat{assumption}{#2Assumption~#1#3}
\crefrangeformat{assumption}{#3Assumptions~#1#4-#5#2#6}
\Crefrangeformat{assumption}{#3Assumptions~#1#4-#5#2#6}

% Automatically resizing commmands
\let\oldexp\exp
\renewcommand{\exp}[1]{\operatorname{exp}\left( #1 \right)}

\newcommand{\Prob}[1]{\mathds{P}\!\!\,\left[#1\right]}
\newcommand{\CProb}[2]{\mathds{P}\!\!\,\left[#1\,\middle|\,#2\right]}
\newcommand{\Exp}[1]{\mathds{E}\!\!\,\left[#1\right]}
\newcommand{\CExp}[2]{\mathds{E}\!\!\,\left[#1\,\middle|\,#2\right]}
\newcommand{\lmax}[1]{\lambda_{\max{}}\!\!\left(#1\right)}
\newcommand{\lmin}[1]{\lambda_{\min{}}\!\!\left(#1\right)}

% Option to manually resize
\newcommand{\expM}[2]{\operatorname{exp}#2( #1 #2)}
\newcommand{\ProbM}[2]{\mathds{P}\!\,#2[#1 #2]}
\newcommand{\CProbM}[3]{\mathds{P}\!\,#3[#1\,#3|\,#2 #3]}
\newcommand{\ExpM}[2]{\mathds{E}\!\,#2[#1 #2]}
\newcommand{\CExpM}[3]{\mathds{E}\!\,#3[#1\,#3|\,#2 #3]}
\newcommand{\lmaxM}[2]{\lambda_{\max{}}\!#2(#1 #2)}
\newcommand{\lminM}[2]{\lambda_{\min{}}\!#2(#1 #2)}

\newcommand{\Cov}[2]{\operatorname{Cov}\!\left(#1, #2\right)}
\newcommand{\CovM}[3]{\operatorname{Cov}\!#3(#1, #2 #3)}
\newcommand{\Var}[1]{\operatorname{Var}\left(#1\right)}
\newcommand{\VarM}[2]{\operatorname{Var}#2(#1 #2)}

% scalars -> normal letters small italics
% vectors -> small roman
% vector of vectors -> small bold roman
% matrics/operators -> large roman
% sets: \mathcal
% sigma-algebras: \mathscr
% probability of a set: \Prob{ ... }


